\section{绪论}
\label{sec:intro}
\subsection{研究背景}
智能座舱中的摄像头可以连续观察驾驶员面部，为表情分析与人机交互提供信息。与普通照片不同，车内画面往往围绕同一驾驶员长时间采集，既包含眉眼、嘴角等表情线索，也包含可以区分个人的五官比例与面部纹理。如何在保留表情信息的同时削弱身份信息，是本项目的出发点。

本项目将隐私处理放在表情识别之前。输入图像先转换为灰度图，再经过人脸检测、对齐与局部匿名化，表情分类器只读取匿名后的灰度人脸。原始图像保留在对照显示一侧，不进入表情推理。这样，匿名化不只是演示界面上的图像效果，而是表情模型输入链路中的实际处理步骤。

\subsection{研究问题}
项目需要同时处理区域控制、信息保留和连续视频三个问题。区域控制要求修改集中在面部内部，保留头发、头部外轮廓与座舱背景；信息保留要求匿名脸仍能表达眼睛开合、眉毛变化和嘴部形态；连续视频则要求相邻帧的处理和判定具有连贯性，能够持续输出表情、置信度与分析速度。

在评价上，匿名化关注本项目自建身份识别模型能否继续匹配原身份，表情识别关注匿名输入与人工标签是否一致，实时性能关注模型完成一帧分析所需的时间。这三个任务分别记录，构成系统的联合评价。

\subsection{主要工作与贡献}
项目以局部残差、几何约束与关键点条件构建统一匿名化模块，在匿名缓存上完成表情训练与校准，并用因果时序处理改善连续输出。Python/PySide6 桌面端提供图片、视频与摄像头分析及模型表现三个页面。训练和运行统一分类输入、标签顺序与归一化方式，使图片实验、离线评价和视频应用共用同一处理约定。

\subsection{文章结构}
第二至第四章依次介绍方案选择、系统设计和训练方法，第五章给出实验结果，第六、七章讨论工程特点并总结。附录列出运行命令、资源、检查与分工。
